% 修士論文（骨組み）
% 組版: platex thesis.tex を2回（相互参照と目次のため）→ dvipdfmx thesis.dvi
% 組版原則は ../principles.md ——プリアンブルへの追加・変更はそちらの規則に従う．
%
% 執筆の親文書は .claude/context.md である．各節の〔執筆メモ〕に対応する節番号を記した．
% 本文の執筆は，残りの測定（特に表現の頑健性，context.md §8.1）が済んでから行う．
\documentclass[report]{jsbook}
\usepackage{amsmath,amssymb}% 本文の数式は context.md の形をそのまま使う（\mathbb・\varnothing に要る）
\usepackage{xurl}% url 互換で，URL をどこでも改行できる（参考文献の Underfull を防ぐ）
\usepackage[dvipdfmx,hidelinks,bookmarksnumbered]{hyperref}
\usepackage{pxjahyper}
\hypersetup{pdftitle={視覚言語モデルにおけるタスク類似性に基づく内部表現整合学習に関する検討},
            pdfauthor={中野友晴}}

% 執筆メモ——書くべき内容と出典（context.md の節）を PDF 上に印字する．
% 本文を書き入れた節から順に，この行を消していく．脱稿時にはすべて消えている．
\newcommand{\draftnote}[1]{%
  \begin{quote}\small〔執筆メモ〕#1\end{quote}}

\title{視覚言語モデルにおけるタスク類似性に基づく\\内部表現整合学習に関する検討}
% 題目は提出時に確定する．候補（杉村先生の指針：疑問形か，明らかにした語で閉じる）：
%   疑問形——タスク類似性に基づく内部表現整合学習により，視覚言語モデルはタスク構造を使うようになるか？
%   閉じ形——視覚言語モデルにおけるタスク類似性に基づく内部表現整合学習がもたらす読み取りやすさと利用の乖離
\author{中野 友晴}

\begin{document}

% ---------------- 表紙 ----------------
\begin{titlepage}
\centering
\vspace*{20mm}
{\Large 修士論文}\\[18mm]
{\LARGE\bfseries 視覚言語モデルにおけるタスク類似性に基づく\\[2mm]内部表現整合学習に関する検討}\\[8mm]
{\small ※題目は提出時に確定（プリアンブルの候補を参照）}\\[28mm]
{\Large 中野 友晴}\\[8mm]
{\large 東京都立大学大学院 システムデザイン研究科\\[1mm]情報科学域}\\[16mm]
{\large 指導教員\quad 杉村 大輔}\\[24mm]
{\large 提出日\quad ※未定}\\ % 提出期限・要提出物は研究科の募集要項で確認する
\vfill
\end{titlepage}

\frontmatter

% ---------------- 概要 ----------------
\chapter{概要}

視覚言語モデル（VLM）は画像と言語を併せて理解し，画像への質問応答など広く応用されている．従来のVLMは出力だけを目標に学習し，答えが同じでも内部表現は入力の揺らぎで大きく動きうる．そこで本研究では，解き方が同じ質問の内部表現を近づける整合学習を検討する．これにより内部表現が解き方に沿って整い，答えの生成に使われると期待できる．実験の結果，正解率をほとんど変えずに解き方を内部表現から読み取りやすくできるが，モデルはそれを使うようにはならず，整合した表現への依存はむしろ下がることがわかった．

\draftnote{上の250字（講演申込と同文）を種に，読み手の知る地点から書き直して800字程度へ広げる（流れは VLM とは何か→従来の課題→提案と期待→結果）．頑健性の測定（context.md §8.1）が済んだら，その結果もここに一文で入る．}

\tableofcontents

\mainmatter

\chapter{序論}
\label{chap:introduction}

\section{背景——出力の頑健性は，表現の頑健性を含意しない}
\label{sec:background}

\draftnote{context.md §1.1．VLM の頑健性は出力で測られてきたこと，Wani ら \cite{wani} が「答えが変わらないまま内部表現が大きく変化する」ことを示したこと（幾何的な加工と文字の重ね書きの違いまで），著者ら自身は対処法を示していないこと．}

\section{本研究の問い}
\label{sec:question}

\draftnote{context.md §1.2．Wani らは同じ入力の揺らぎを見ており，別々の質問どうしの近さは問われていない——そこから入る．中心の問いは一つ：タスク構造（署名）を内部表現に明示的に整合させたとき，内部表現と答えの生成に何が起きるか．正解率の向上は目的ではないことをここで明言する．}

\section{貢献}
\label{sec:contribution}

\draftnote{context.md §2．貢献は「読み取れること」と「使われること」を分けて示したこと．Elazar ら \cite{elazar} の「除く」方向に対し，本研究は損失で「足す」方向から構成的に示す．VLsI \cite{vlsi} は傍証にとどめ，中核の支柱には据えない（教師あり蒸留と自己整合の違い）．}

\section{本論文の構成}
\label{sec:organization}

\draftnote{各章の一文紹介．最後に書く．}

\chapter{関連研究}
\label{chap:related-work}

% context.md に独立した関連研究章は無い．references/ の検証済みノート（1論文1ファイル）を親文書として書く．
% 本研究との距離（何を借り，どこが違うか）を各節の最後に必ず一文で置く．

\section{VLM の内部表現の不安定性}
\label{sec:rw-instability}

\draftnote{references/hidden-instability.md．Wani ら \cite{wani} の設定（言語側最終層，Qwen3-VL・LLaVA-OneVision，SEEDBench・MMMU・POPE）と指標（無関係画像対の距離分布を基準にした比）．本研究が借りるのは指標の形だけである．}

\section{プローブ——読み取れることと使われることの区別}
\label{sec:rw-probing}

\draftnote{references/probing-classifiers.md・amnesic-probing.md・task-format.md．Belinkov \cite{belinkov} の総覧，Elazar ら \cite{elazar} の amnesic probing（除いて測る・Rand 対照），Sahoo ら \cite{sahoo} の「プローブは task format を検出しうる」という警告．本研究の測定設計（言い回しプローブ・書き出し分離・残差化）がどの批判に応えるものかをここで対応づける．}

\section{対照学習}
\label{sec:rw-contrastive}

\draftnote{references/supcon.md．SupCon \cite{supcon} の損失と射影 head．本研究は射影を挟まず（反証可能性のため）生の表現に掛ける点が異なる——詳細は\ref{sec:where}節．}

\section{内部表現に手を入れる先行例}
\label{sec:rw-representation}

\draftnote{references/vlsi.md．VLsI \cite{vlsi}（層ごとの特徴を自然言語化して蒸留，正解率向上）．外部教師の有無・表現を直接近づけるか否かの違い．}

\section{GQA と functional program}
\label{sec:rw-gqa}

\draftnote{references/gqa.md．GQA \cite{gqa} の構成と functional program．program が付いていることが採用理由であること．}

\chapter{提案手法}
\label{chap:method}

\section{タスク署名——解き方を，比べられる形にする}
\label{sec:signature}

\draftnote{context.md §3.1．functional program の例（``Who is wearing the dress?'' の3段の表），署名の定義，引数を捨てる理由，署名と質問文の交絡（自動生成ゆえの相関，演算子名が質問文に現れる10種）．}

\section{整合損失}
\label{sec:loss}

\draftnote{context.md §3.2．SupCon 式（$\mathcal{L}_{\mathrm{align}}$）と各項の読み，$\mathcal{I}$ から外れる質問が本研究のバッチ構成では生じないこと，温度 $\tau$（学習可能，初期値0.07），全体の目的関数 $\mathcal{L} = \mathcal{L}_{\mathrm{LM}} + \lambda\mathcal{L}_{\mathrm{align}}$（$\lambda=0.1$）．}

\section{どこに掛けるか——射影を挟まない中間層の一点}
\label{sec:where}

\draftnote{context.md §3.3．第16層・質問を読み終えた位置（$H=960$）の選定理由（causal attention のもとで画像と質問文を読み終えている唯一の位置），射影 head を挟まない理由（反証可能性），それでも構造が使われない余地が残ること（多対一の写像）——この余地を測るのが\ref{sec:usage}節である．}

\chapter{実験設定}
\label{chap:experiments}

\section{モデルとデータ}
\label{sec:model-data}

\draftnote{context.md §4.1．SmolVLM-500M-Instruct \cite{smolvlm}（32層，$H=960$），GQA \texttt{train\_balanced} 1エポック（884問除外，942,116問），バッチ構成（16署名 × 2件），学習率・勾配打ち切り，9回の学習（3群 × 3乱数種，1回約2時間53分），\texttt{testdev\_balanced} 12,578問で報告する理由と3分割の排他性．}

\section{三つの群}
\label{sec:groups}

\draftnote{context.md §4.2．baseline／proposal／ablation の定義．ablation は「ラベルの意味だけを除く」群であり表現を壊す群ではないこと，この対照が答えるのは「ラベルに意味が要るか」までで「署名である必要があるか」ではないこと（→\ref{sec:future}節）．}

\section{乱数種3つの，まとめ方}
\label{sec:statistics}

\draftnote{context.md §4.3．対応のある $t$ 検定と95\%信頼区間（$n=3$，$t_{2,0.975}=4.303$），「区間が0を含む＝差が無い，ではない」という読み方，baseline 3本で測る乱数種由来のばらつき（幅1.14pt）．}

\section{整合が働いたことの確認}
\label{sec:sanity}

\draftnote{context.md §4.4．分離度 $\delta$（Cohen's d）の定義と3群の値（0.816／2.378／0.451），proposal は「同署名を近づけた」のではなく異署名がそれ以上に離れた結果であること，ablation の表現が一方向に潰れたことと $\log 31$ への一致（ラベルに情報が無いことの裏づけ）．}

\chapter{結果}
\label{chap:results}

% 主結果は三つ．順序も見出しの立て方も context.md §5 と同じにする．
% 測定方法の詳細（プローブの当てはめ・三段の測定）は付録\ref{chap:appendix-verification}に置き，
% 本章は付録を読まなくても筋が通るように書く——context.md の本文／付録A の分担をそのまま保つ．

\section{署名の読み取りやすさは，大きく上がる}
\label{sec:decodability}

\draftnote{context.md §5.1（当てはめの詳細は付録\ref{sec:app-probe}）．読み取り精度の表（82.9\% → 95.9\%，+13.0pt；ablation 81.0\%；言い回しプローブ 69.7\%；床 28.8\%；書き出し分離で +14.5pt）．「上がったこと自体は発見ではない」——確定させるのは増分の帰属であるという書き方を保つ．整合前の baseline の読み（幾何は近くないが線形には読める）もここ．組成への波及は付録\ref{sec:app-composition}へ送る．}

\section{正解率は，ほとんど動かない}
\label{sec:accuracy}

\draftnote{context.md §5.2．乱数種ごとの正解率の表，proposal−baseline の区間 $[-1.81, +1.70]$ の計算過程，「差が無い」ではなく上限を言う読み方．上がらない理由の層別は付録\ref{sec:app-frequency}へ送る．}

\section{読み取れるようになった構造は，使われていない}
\label{sec:usage}

\draftnote{context.md §5.3（測定方法は付録\ref{sec:app-elasticity}）．弾性の表（3幅で −28.0／−21.6／−18.0\%，ablation は下がらない），特権の消失（+2.82pt → 差なし），一次近似であることの留保（「まったく使っていない」ではなく「頼る度合いが下がった」）．}

\section{表現の頑健性——整合した表現は，入力の揺らぎに安定するか}
\label{sec:robustness}

\draftnote{【未測定——この測定が済むまで本論文は書き始めない】context.md §8.1．出発点（\ref{sec:background}節）に応える測定：3群に正解を変えない画像加工を施し，第16層の表現の変化量を無関係画像対の距離分布を基準にした比で比べる．結果が3通りのどれか（安定性を高める／形式だけ／脆くなる）で\ref{chap:conclusion}章の枠組みが変わる．9回の学習の重みは E:\textbackslash runs（USB バックアップあり）にある．}

\chapter{結論}
\label{chap:conclusion}

\section{結論と貢献}
\label{sec:conclusion}

\draftnote{context.md §6．三つの主結果を一文で束ねる．貢献（足す側から構成的に示した），「この研究から言えないこと」の列挙をここでも省かない——頑健性（\ref{sec:robustness}節）の結果が出たら，束ねる一文はそれを含む形に書き直す．}

\section{限界}
\label{sec:limitations}

\draftnote{context.md §7．設計の一点性（$\lambda=0.1$・第16層・単一モデル・単一データセット・1エポック），署名が解き方の正しい形式化かは未検証（GQA 公式 consistency との食い違い52.1\%の読み方を含む），「使われていない」は第16層・一次近似の範囲であること．}

\section{今後の課題}
\label{sec:future}

\draftnote{context.md §8．執筆時点で残っているものだけを書く（済んだ測定は\ref{chap:results}章へ移す）：稀な署名の表現側の層別（§8.2），答えの型ラベルでの整合——主語の確定（§8.3），$\lambda$ と層を振る・別モデルと別データセット（§8.4，署名に代わる形式化が要ること）．}


% ---------------- 謝辞 ----------------
\chapter*{謝辞}
\addcontentsline{toc}{chapter}{謝辞}

\draftnote{指導教員・研究室・（あれば）計算資源等への謝辞．最後に書く．}

% ---------------- 参考文献 ----------------
% context.md の参考文献と references/ のノートに対応．並びは本文での初出順（../principles.md）．
\begin{thebibliography}{99}
\bibitem{wani} F. A. Wani et al.: ``Same Answer, Different Representations: Hidden Instability in VLMs,'' arXiv:2602.06652, 2026.
\bibitem{elazar} Y. Elazar et al.: ``Amnesic Probing: Behavioral Explanation with Amnesic Counterfactuals,'' TACL, vol. 9, pp. 160--175, 2021.
\bibitem{vlsi} B.-K. Lee et al.: ``VLsI: Verbalized Layers-to-Interactions from Large to Small Vision Language Models,'' CVPR, 2025.
\bibitem{belinkov} Y. Belinkov: ``Probing Classifiers: Promises, Shortcomings, and Advances,'' Computational Linguistics, vol. 48, no. 1, pp. 207--219, 2022.
\bibitem{sahoo} S. Sahoo et al.: ``Linear Probes Detect Task Format, Not Reasoning Mode in Language Model Hidden States,'' 6th Workshop on Trustworthy NLP, ACL 2026.
\bibitem{supcon} P. Khosla et al.: ``Supervised Contrastive Learning,'' NeurIPS, 2020.
\bibitem{gqa} D. A. Hudson and C. D. Manning: ``GQA: A New Dataset for Real-World Visual Reasoning and Compositional Question Answering,'' CVPR, 2019.
\bibitem{smolvlm} SmolVLM-500M-Instruct, \url{https://huggingface.co/HuggingFaceTB/SmolVLM-500M-Instruct}
\bibitem{hydra} T. McGrath et al.: ``The Hydra Effect: Emergent Self-repair in Language Model Computations,'' arXiv:2307.15771, 2023.
\bibitem{workspace} W. Gurnee et al.: ``Verbalizable Representations Form a Global Workspace in Language Models,'' Transformer Circuits Thread, 2026.
\end{thebibliography}

% ---------------- 付録 ----------------
\appendix
\chapter{補助的な検証}
\label{chap:appendix-verification}

% context.md 付録A をそのまま写す．本文（第\ref{chap:results}章）はこの付録を読まなくても筋が通ること．
% 審査で「方法の詳細は本文に」と求められたら，A.1 と A.4 は本文の対応節へ昇格させる（その判断は指導教員と）．

\section{「読み取りやすさ」の測定と，二つの対照}
\label{sec:app-probe}

\draftnote{context.md §A.1．プローブの定義・当てはめと打ち切り（10問以上の署名に限る，1/2分割，early stopping）・床の定義，言い回しプローブと書き出し分離の作り方，ablation の81.0\%の読み方，残差化（Sahoo ら \cite{sahoo} への応答，形式3変量で0.35ptしか動かない）——残差化はこの付録限定とする．}

\section{演算子の組成も，読み取れるようになる}
\label{sec:app-composition}

\draftnote{context.md §A.2．「学習が教えたことを読み返しているだけ」への反論の設計（署名単位の分割・演算子集合ラベル），AUC と完全一致率の表（36.3\% → 44.5\%），言い回しプローブとの比較の読み方の注意．}

\section{正解率が上がらない理由——頻度層別と，正誤で分けたプローブ}
\label{sec:app-frequency}

\draftnote{context.md §A.3．署名の学習出現回数による5層の層別表（利得の現れる層は無い，1--10層の −2.50pt は区間が0を跨ぐ），正誤で分けたプローブ（誤答側でも96.4\%——誤りは署名より後の段階），proposal に掛かる留保．稀な署名の表現側の測定（§8.2）が済んだらここに足す．}

\section{「使われていない」を，どう測ったか——三段の測定}
\label{sec:app-elasticity}

\draftnote{context.md §A.4．中立／不一致の二択，自己修復 \cite{hydra} ゆえに壊す方法を採らないこと，三段の設計（$U_m$・J-lens \cite{workspace} による引き戻し $J(U_m)$・弾性 $e_p(u)$）と判定量 $\Delta_m$／$\Delta'_m$，結果の表と各段の読み（第1段は決着に使わない，など），64問で足りることの確認．}

\section{整合する層を変えた試走——乱数種1つ}
\label{sec:app-layers}

\draftnote{context.md §A.5．第8層・第24層（乱数種42）の表，読み取り精度は層を変えても同じ高さに着くこと，区間を作れない旨．乱数種を揃える追試（§8.4）が済んだら本文へ昇格を検討する．}


\end{document}
